\subsection{Coproduct for symmetric tensors}

Let M be a free A-module, and \(\Delta_{M} = \Delta\) the diagonal homomorphism \(x \mapsto (x, x)\) of M into \(M \oplus M\) . Let \(c_{M} = c\) be the unital homomorphism of graded A-algebras composed of the homomorphisms:

\[
\mathbf {T S} (\mathrm{M}) \xrightarrow {\mathbf {T S} (\Delta)} \mathbf {T S} (\mathrm{M} \oplus \mathrm{M}) \xrightarrow {\sigma} \mathbf {T S} (\mathrm{M}) \otimes \mathbf {T S} (\mathrm{M})
\]

where \(\sigma\) is the canonical isomorphism. Equipped with \(c\) , \(\mathbf{TS}(\mathbf{M})\) is a graded A-cogebra.

For all \(x \in \mathbf{M}\) and each integer \(p \geqslant 0\) we have \(\mathbf{TS}(\Delta) (\gamma_p(x)) = \gamma_p((x, x))\) , hence by (8),

\[
c \left(\gamma_ {p} (x)\right) = \sum_ {r + s = p} \gamma_ {r} (x) \otimes \gamma_ {s} (x).\tag{10}
\]

In particular

\[
c (x) = x \otimes 1 + 1 \otimes x.\tag{11}
\]

Let \((x_{i})_{i\in I}\) be a family of elements of \(\mathbf{M}\) , and for \(\nu \in \mathbf{N}^{(I)}\) put \(x_{\nu} = \prod_{i\in I}\gamma_{\nu_i}(x_i)\) . Then

\[
c \left(x _ {\nu}\right) = \sum_ {\rho + \sigma = \nu} x _ {\rho} \otimes x _ {\sigma}.\tag{12}
\]

This follows from (10) because \(c\) is an algebra homomorphism.

PROPOSITION 8. --- Let M be a free A-module, then with its algebra and cogebra structures, TS(M) is a graded commutative and cocommutative bigebra. The counit is the A-linear mapping \(\varepsilon: \mathbf{TS}(\mathbf{M}) \to \mathbf{TS}^0(\mathbf{M}) = \mathbf{A}\) which is zero on \(\mathbf{TS}^p(\mathbf{M})\) for \(p > 0\) and such that \(\varepsilon(1) = 1\) .

We know that the A-algebra TS(M) is associative, commutative and unital. On the other hand, the coproduct is by construction a homomorphism of graded algebras; now the fact that the cogebra TS(M) is coassociative and cocommutative follows by an easy calculation from the Formula (10). The mapping \(\varepsilon\) of TS(M) into A is a homomorphism of graded algebras such that \(\varepsilon(1)=1\) . Finally for every \(x\in M\) we have \(\varepsilon(\gamma_{p}(x))=0\) if p\textgreater0, \(\varepsilon(\gamma_{0}(x))=1\) ; if we bear in mind (10), this shows that \((\varepsilon\otimes1)\circ c=(1\otimes\varepsilon)\circ c=\mathrm{Id}_{\mathbf{TS}(\mathbf{M})}\) ; thus \(\varepsilon\) is the counit of TS(M).

PROPOSITION 9. --- Let M and N be free A-modules and u an A-homomorphism of M into N; then TS(u) is a bigebra homomorphism.

For we have \(\Delta_{\mathbf{N}} \circ u = (u, u) \circ \Delta_{\mathbf{M}}\) , hence the diagram

\[
\begin{array}{c} \mathsf {T S (M)} \xrightarrow {\mathsf {T S} (\Delta_ {M})} \mathsf {T S (M \oplus M)} \xrightarrow {\sigma} \mathsf {T S (M)} \otimes \mathsf {T S (M)} \\ \mathsf {T S (u)} \Bigg | _ {\mathsf {T S (u)}} \\ \mathsf {T S (N)} \xrightarrow {\mathsf {T S} (\Delta_ {N})} \mathsf {T S (N \oplus N)} \xrightarrow {\tau} \mathsf {T S (N)} \otimes \mathsf {T S (N)}, \end{array}
\]

where \(\sigma\) and \(\tau\) are canonical isomorphisms, is commutative (IV, p.~49). Thus \(c_{\mathrm{N}} \circ \mathbf{T}\mathbf{S}(u) = (\mathbf{T}\mathbf{S}(u) \otimes \mathbf{T}\mathbf{S}(u)) \circ c_{\mathrm{M}}\) .

PROPOSITION 10. --- Let M be a free A-module, then the primitive elements (III, p.~602) of the bigebra TS(M) are the elements of M.

Let \((e_i)_{i\in I}\) be a basis of \(\mathbf{M}\) , and for \(\nu \in \mathbf{N}^{(I)}\) put \(e_{\nu} = \prod_{i\in I}\gamma_{\nu_i}(e_i)\) . Let

\(z = \sum_{\nu \in \mathbf{N}^{(I)}} \lambda_{\nu} e_{\nu}\) be an element of \(\mathbf{T}\mathbf{S}(\mathbf{M})\) , then by (12) we have

\[
c (z) = \sum_ {\nu} \lambda_ {\nu} \sum_ {\rho , \sigma \in \mathbf {N} ^ {(I)}, \rho + \sigma = \nu} e _ {\rho} \otimes e _ {\sigma} = \sum_ {\rho , \sigma} \lambda_ {\rho + \sigma} e _ {\rho} \otimes e _ {\sigma}
\]

hence

\[
c (z) - 1 \otimes z - z \otimes 1 = \sum_ {\rho \neq 0, \sigma \neq 0} \lambda_ {\rho + \sigma} e _ {\rho} \otimes e _ {\sigma} - \lambda_ {0} e _ {0} \otimes e _ {0}.
\]

and so

\[
\begin{array}{r l} z \text {   primitive   } & \Leftrightarrow \lambda_ {\rho + \sigma} = 0 \text {   when   } \rho \neq 0 \text {   and   } \sigma \neq 0 \text {   and   } \lambda_ {0} = 0 \\ & \Leftrightarrow \lambda_ {\nu} = 0 \text {   when   } | \nu | \neq 1 \\ & \Leftrightarrow z \in M. \end{array}
\]

